% ==============================================================================
% SECCIÓN 1: IMAGEN DE EMPRESA
% ==============================================================================
\section{Imagen de Empresa}

\subsection{Identidad Corporativa: CLENCH Software Development}
El equipo de trabajo opera bajo la identidad corporativa y marca registrada \textbf{CLENCH Software Development}. La iniciativa surge con un marcado carácter de \textbf{Aprendizaje-Servicio}, aunando la excelencia en ingeniería de software con la responsabilidad social y la inclusión de colectivos vulnerables.

\begin{center}
  \includegraphics[width=6.0cm]{images/logo_clench.png}
\end{center}

\begin{itemize}[leftmargin=*]
  \item \textbf{Denominación Oficial:} CLENCH Software Development.
  \item \textbf{Lema y Propósito:} \textit{«Tecnología accesible para transformar vidas y potenciar la autonomía personal»}.
  \item \textbf{Misión:} Diseñar, implementar y validar soluciones digitales accesibles, robustas e intuitivas que eliminen barreras cognitivas y motrices en la rutina diaria de estudiantes con Necesidades Específicas de Apoyo Educativo (NEAE).
  \item \textbf{Visión:} Consolidar un modelo de equipo multidisciplinar altamente cohesionado, referente en buenas prácticas de metodologías ágiles y cumplimiento estricto de las directrices universales de accesibilidad (WCAG~2.1~AA).
  \item \textbf{Valores Corporativos:}
  \begin{enumerate}
    \item \textbf{Accesibilidad Universal y Empatía:} La persona usuaria (estudiantes PTVAL y profesorado de la Fundación Purísima Concepción) es el núcleo de cada decisión de diseño.
    \item \textbf{Rigor Técnico y Calidad Continua:} Compromiso innegociable con el código limpio, análisis estático, cobertura de pruebas y estabilidad.
    \item \textbf{Transparencia y Coordinación:} Comunicación fluida, reparto equitativo de esfuerzo y sinceridad en la estimación de plazos y riesgos.
    \item \textbf{Ética y Protección de Datos:} Respeto absoluto a la privacidad, dignidad y normativa RGPD en el tratamiento de información asistencial.
  \end{enumerate}
\end{itemize}
